\section*{Capítulo \ref{ch:b1:lab-techniques-1} --- Técnicas de laboratorio I: seguridad, medida y separación}

\begin{solution}{exo:b1:lab-techniques-1:1}
\emph{Peligro}: la exige H225 (líquido muy inflamable, categoría 2), y la
etiqueta lleva la palabra más grave. H225 es un \omterm{def:b1:lab-techniques-1:hazard}{peligro} físico; H319
(irritación ocular grave) y H336 (somnolencia o vértigo) son \omterm{def:b1:lab-techniques-1:hazard}{peligros}
para la salud. Precauciones: ninguna llama ni chispa cerca, una botella
cerrada; gafas; trabajar en un lugar ventilado o en una vitrina de gases.
% ledger: ghs:acetone, hstat:H225, hstat:H319, hstat:H336
\end{solution}

\begin{solution}{exo:b1:lab-techniques-1:2}
(a) Un \omterm{def:b1:lab-techniques-1:hazard}{peligro}, una propiedad del ácido. (b) Un \omterm{def:b1:lab-techniques-1:hazard}{riesgo}, pequeño porque
la exposición es pequeña. (c) Un \omterm{def:b1:lab-techniques-1:hazard}{peligro}. (d) Un \omterm{def:b1:lab-techniques-1:hazard}{riesgo}: el mismo
\omterm{def:b1:lab-techniques-1:hazard}{peligro}, con mayor exposición (vapor, salpicaduras) en caliente y al
aire.
\end{solution}

\begin{solution}{exo:b1:lab-techniques-1:3}
$R_f = 2.1/6.0 = 0.35$ y $4.2/6.0 = 0.70$. La segunda mancha tiene el
mismo $R_f$ que la referencia: la muestra puede contenerla (es
coherente), pero un mismo $R_f$ con un solo eluyente no prueba la
identidad; un segundo eluyente, o un depósito conjunto de muestra y
referencia que siga dando una sola mancha, refuerza la conclusión. La
mancha a 0.35 es con certeza otra especie.
\end{solution}

\begin{solution}{exo:b1:lab-techniques-1:4}
Semianchura de \qty{0.05}{mL}, rectangular: $u = 0.05/\sqrt3 = \qty{0.029}{mL}$.
Para una diferencia de dos lecturas independientes,
$u = \sqrt{2} \times 0.029 = \qty{0.041}{mL}$.
\end{solution}

\begin{solution}{exo:b1:lab-techniques-1:5}
Media, \qty{0.25100}{g}; desviaciones $+2$, $-2$, $+5$, 0, $-5$ (en
\qty{e-4}{g}), suma de cuadrados \qty{58e-8}{g^2}, $s = \sqrt{58 \times
10^{-8}/4} = \qty{3.8e-4}{g}$; $u = s/\sqrt5 = \qty{1.7e-4}{g}$;
$U = \qty{3.4e-4}{g}$: $m = (0.25100 \pm 0.00034)$~g, $k = 2$.
\end{solution}

\begin{solution}{exo:b1:lab-techniques-1:6}
Una sola porción: $f = 1/(1 + 3 \times 30/50) = 0.357$; se extrae el
\qty{64.3}{\percent}. Dos de \qty{15}{mL}: $f = 1.9^{-2} = 0.277$, el
\qty{72.3}{\percent}. Tres de \qty{10}{mL}: $f = 1.6^{-3} = 0.244$, el
\qty{75.6}{\percent}.
\end{solution}

\begin{solution}{exo:b1:lab-techniques-1:7}
$z = 0.0012/\sqrt{0.0006^2 + 0.0002^2} = 0.0012/0.00063 = 1.9 \le 2$:
compatibles, por poco. Con \qty{0.0004}{mol/L}: $z = 0.0012/0.00045 = 2.7$,
no compatibles: la \omterm{def:b1:titration-methods:titration}{valoración} más precisa revela una diferencia real (un
error en la preparación, o en la \omterm{def:b1:titration-methods:titration}{valoración}).
\end{solution}

\begin{solution}{exo:b1:lab-techniques-1:8}
Q: poco soluble en frío, muy soluble en caliente. P apenas disuelve el
sólido mejor en caliente que en frío (nada cristaliza al enfriar); R
mantiene demasiado disuelto en frío (grandes pérdidas). Con Q,
\qty{5.0}{g} necesitan al menos $5.0/12 \times 100 = \qty{42}{mL}$ de
disolvente en ebullición, que en frío mantienen disueltos
$0.3 \times 0.417 = \qty{0.13}{g}$: pueden recuperarse como mucho
\qty{4.9}{g}, el \qty{97.5}{\percent} (antes de cualquier pérdida en el
filtro).
\end{solution}

\begin{solution}{exo:b1:lab-techniques-1:9}
Etanol ($n_D = 1.3611$ a \qty{20}{\celsius}); el agua (1.333) y el
ciclohexano (1.4266) quedan excluidos. El índice varía con la
temperatura, de modo que un valor no significa nada sin ella. Un valor
de 1.350, entre el agua y el etanol, sugiere una mezcla, como un etanol
que contiene agua.
% ledger: nD:water, nD:ethanol, nD:cyclohexane
\end{solution}

\begin{solution}{exo:b1:lab-techniques-1:10}
Con $t = KV/V_a$ y $n$ porciones, $\ln f_n = -n\ln(1 + t/n)$; cuando
$n \to \infty$, $n \ln(1 + t/n) \to t$ (pues $\ln(1 + \varepsilon) \sim
\varepsilon$), y $f_n$ decrece hacia $\mathrm{e}^{-t}$
(\cref{prop:b1:lab-techniques-1:multiple-extractions}). Aquí
$t = 3 \times 30/50 = 1.8$, $\mathrm{e}^{-1.8} = 0.165$: como mucho, el
\qty{83.5}{\percent}. Tres porciones ya dan el \qty{75.6}{\percent}, el
\qty{91}{\percent} del límite; alcanzar el \qty{99}{\percent} de este
exige 32 porciones de menos de \qty{1}{mL}: más allá de dos o tres
porciones, la ganancia no compensa el trabajo.
\end{solution}

\begin{solution}{exo:b1:lab-techniques-1:11}
Con $t = x - x_0$: $\int_{-a}^{a} (a - |t|)/a^2 \,\mathrm{d}t =
2 \int_0^a (a - t)/a^2 \,\mathrm{d}t = 2 (a^2/2)/a^2 = 1$. La media es
$x_0$ por simetría; la varianza es $2 \int_0^a t^2 (a - t)/a^2
\,\mathrm{d}t = \frac{2}{a^2}\left(\frac{a^4}{3} - \frac{a^4}{4}\right) =
\frac{a^2}{6}$, de modo que $u = a/\sqrt6 = 0.41\,a$, frente a
$a/\sqrt3 = 0.58\,a$ en el caso rectangular: saber que el centro es más
probable reduce la incertidumbre.
\end{solution}

\begin{solution}{exo:b1:lab-techniques-1:12}
$c = m/(MV) = 0.25100/(204.2 \times 0.10000) = \qty{0.012292}{mol/L}$.
\omterm{def:b1:lab-techniques-1:uncertainty}{Incertidumbres relativas}: masa, $1.7 \times 10^{-4}/0.2510 = 6.8 \times
10^{-4}$; volumen, $0.06/100.00 = 6.0 \times 10^{-4}$; combinada,
$\sqrt{(6.8^2 + 6.0^2)} \times 10^{-4} = 9.1 \times 10^{-4}$, es decir,
\qty{0.091}{\percent}; $u(c) = \qty{1.1e-5}{mol/L}$, $U = \qty{2.2e-5}{mol/L}$:
$c = (0.012292 \pm 0.000022)$~mol/L, $k = 2$. Domina la pesada, por
poco; las dos fuentes son del mismo orden.
\end{solution}

\begin{solution}{pb:b1:lab-techniques-1:1}
\textbf{1.} GHS02, inflamable; GHS05, corrosivo; GHS06, toxicidad aguda;
GHS07, nocivo o irritante.
\textbf{2.} Ácido salicílico: \emph{Peligro}, por H318 (lesiones oculares
graves, categoría 1). Aspirina: \emph{Atención} (solo H302).
\textbf{3.} El \omterm{def:b1:lab-techniques-1:hazard}{peligro} es fijo, la exposición no: un volumen pequeño,
medido y usado en una vitrina de gases, con guantes, gafas y bata, la
botella cerrada enseguida, ninguna llama cerca.
\textbf{4.} H314 (quemaduras graves en la piel y lesiones oculares
graves): gafas y guantes. H332 (nocivo en caso de inhalación) y H226
(líquidos y vapores inflamables): la vitrina de gases y la ausencia de
llamas.
\textbf{5.} Aclarar cuidadosamente con agua durante varios minutos,
quitando las lentes de contacto si es posible, y seguir aclarando
(P305+P351+P338); después, consultar a un médico.
\textbf{6.} Al recipiente de residuos acuosos, tras neutralizarlo, no al
fregadero.
\textbf{7.} \ce{C7H6O3 + C4H6O3 -> C9H8O4 + C2H4O2}: C $7 + 4 = 9 + 2$,
H $6 + 6 = 8 + 4$, O $3 + 3 = 4 + 2$.
\textbf{8.} $2.00/138.0 = \qty{1.449e-2}{mol}$; como mucho,
$\num{1.449e-2} \times 180.0 = \qty{2.61}{g}$ de aspirina.
\textbf{9.} El sólido bruto contiene ácido salicílico sin reaccionar,
ácido etanoico y agua: su masa no es la de la aspirina, y su pureza es
desconocida.
\textbf{10.} Lo que queda disuelto tras enfriar se pierde: una
\omterm{def:b1:precipitation:solubility}{solubilidad} pequeña en frío y grande en caliente significa que el
producto se disuelve en caliente y vuelve casi por completo en frío.
\textbf{11.} Cada mililitro de más mantiene disuelta en frío su parte de
producto.
\textbf{12.} Un crecimiento lento da \omterm{def:b1:crystals-metals:crystal}{cristales} más grandes y regulares,
que atrapan menos aguas madres e impurezas; el hielo rebaja la
\omterm{def:b1:precipitation:solubility}{solubilidad} y, por tanto, la pérdida.
\textbf{13.} $\qty{0.045}{L} \times \qty{3.3}{g/L} = \qty{0.15}{g}$.
\textbf{14.} $1.95/2.61 = \qty{75}{\percent}$ (\qty{74.7}{\percent}).
\textbf{15.} $807/6 = \qty{134.5}{\celsius}$.
\textbf{16.} Desviaciones $-0.5$, $0.5$, $-1.5$, $0.5$, $-0.5$, $1.5$; suma
de cuadrados, 5.5; $s = \sqrt{5.5/5} = \qty{1.05}{\celsius}$.
\textbf{17.} $u_A = 1.05/\sqrt6 = \qty{0.43}{\celsius}$.
\textbf{18.} Semianchura de \qty{0.5}{\celsius}: $0.5/\sqrt3 = \qty{0.29}{\celsius}$.
\textbf{19.} $1/\sqrt3 = \qty{0.58}{\celsius}$;
$u = \sqrt{0.428^2 + 0.289^2 + 0.577^2} = \sqrt{0.600} = \qty{0.77}{\celsius}$.
\textbf{20.} $U = 2 \times 0.775 = \qty{1.5}{\celsius}$:
$T = (134.5 \pm 1.5)$~\unit{\celsius}, $k = 2$.
\textbf{21.} Funde bajo y en un intervalo de \qty{8}{\celsius}: es impuro
(ácido salicílico, ácido etanoico, agua).
\textbf{22.} Dado a la unidad, el valor está a menos de
$\pm\qty{0.5}{\celsius}$: $u_{\mathrm{ref}} = 0.5/\sqrt3 = \qty{0.29}{\celsius}$.
\textbf{23.} $2.3/6.0 = 0.38$ y $3.3/6.0 = 0.55$. El producto bruto
contenía ácido salicílico (mismo $R_f$ que la referencia) y aspirina; tras
la \omterm{def:b1:lab-techniques-1:recrystallisation}{recristalización} solo queda la mancha de la aspirina: el ácido
salicílico se ha eliminado, al menos por debajo de lo que detecta la
placa.
\textbf{24.} La aspirina se descompone al fundir; calentada lentamente,
tiene tiempo de descomponerse, y la mezcla formada funde más abajo. El
valor tabulado se refiere a un calentamiento rápido, como en el banco.
\textbf{25.} $z = |134.5 - 135|/\sqrt{0.600 + 0.083} = 0.5/0.827 =$
\textbf{0.60}: muy por debajo de 2; los puntos de fusión medido y
tabulado son compatibles.
\end{solution}
